\chapter{Introduction}
\label{chap:introduction}

Organizational psychology has established that leadership style is a primary determinant of team climate and productivity in human groups~\cite{mathieu2017teamwork, goleman2000leadership}.
However, whether these behavioral patterns also apply to artificial agent teams remains an open question~\cite{muralidharan2025lessonshumanteamsapplied, kwak2026leadershipcoordinationcontrolbehavioral}.
This thesis investigates this by applying Goleman's six leadership styles as distinct system-prompt conditions for a Boss agent in a heterogeneous \ac{LLM} multi-agent architecture, treating supervisory framing as an experimental variable rather than a fixed architectural choice~\cite{guo2024embodiedllmagentslearn, faulkner2026evaluatingcooperationllmsocial}.

\section{Motivation}
\label{sec:motivation}

As \ac{LLM} architectures evolve from single-agent setups toward complex \acp{MAS}, understanding effective orchestration has become a critical challenge~\cite{guo2024largelanguagemodelbased, wang2026multiagentsystemsclassicalparadigms}. 
In heterogeneous \ac{MAS} environments, specialized autonomous agents, such as coders, writers, and reviewers, are assigned highly cooperative roles to execute complex software and data pipelines~\cite{qian2024chatdevcommunicativeagentssoftware, hong2024metagptmetaprogrammingmultiagent}.
However, organizing these expert agents into productive collectives requires managing conversational friction, coordination bottlenecks, and consensus integration~\cite{wu2023autogenenablingnextgenllm, muralidharan2025lessonshumanteamsapplied}.
While existing research primarily focuses on algorithmic task-decomposition mechanisms~\cite{liu2025selectthendecomposeempiricalanalysisadaptive} or structural communication topologies~\cite{chen2023agentversefacilitatingmultiagentcollaboration}, 
human organizational management offers a rich framework of leadership styles that remains largely unexplored in AI orchestration~\cite{mathieu2017teamwork, barrick1998personalityteamwork}. 
In human teams, a leader's managerial approach directly impacts both objective task performance and subjective team dynamics~\cite{mathieu2017teamwork, ayman1995contingency, goleman2000leadership}. 
Grounded in Daniel Goleman's classical framework of emotional intelligence in leadership, managers adopt distinct styles (Coercive, Authoritative, Affiliative, Democratic, Pacesetting, and Coaching) to guide group behavior~\cite{goleman2000leadership}.

Translating these human paradigms to artificial agents presents a compelling question for \ac{MAS} design: 
because \acp{LLM} are trained on vast human interactions, a supervisor's managerial framing may systematically alter how worker agents coordinate and resolve task friction~\cite{park2023generativeagentsinteractivesimulacra, salewski2023incontextimpersonationrevealslarge, shanahan2023roleplaylargelanguagemodels, serapiogarcia2025personalitytraitslargelanguage}. 
This thesis systematically evaluates Goleman's leadership styles within a multi-agent framework to determine their direct impact on task performance and work satisfaction. 
Specifically, this study investigates whether established human management strategies translate into higher task accuracy and efficiency in AI-to-AI collaboration, or if the unique characteristics of artificial agents demand an entirely adapted approach to AI leadership.

To investigate this, an empirical \ac{MAS} was developed in which a supervisor agent guides a team of specialized worker agents through collaborative data science tasks of varying complexity~\cite{dang2025multiagentcollaborationevolvingorchestration, muralidharan2025lessonshumanteamsapplied, wang2025talkstructurallyacthierarchically}.
By holding worker prompts and tasks constant while supervisor prompts are systematically varied across Goleman's six leadership styles and a Baseline, the resulting effects on objective deliverable quality, operational cost, and subjective team dynamics are measured in this study~\cite{goleman2000leadership, kwak2026leadershipcoordinationcontrolbehavioral, muralidharan2025lessonshumanteamsapplied}.

\section{Research Question}
\label{sec:research-question}

To investigate how managerial framing affects AI-to-AI orchestration, this thesis addresses the following primary research question:

\textit{How do Goleman's six leadership styles affect output quality, operational efficiency, communication sentiment, and worker satisfaction in a heterogeneous multi-agent system?}

To answer this question, this study systematically varies the managerial framing of an \ac{LLM} supervisor agent leading a team of specialized worker agents across data science tasks of varying structural complexity. 
The impact of these leadership strategies is evaluated across two core dimensions:

\begin{itemize}
    \item \textbf{Objective Task Performance:} Evaluating how leadership styles influence technical execution, output quality, and computational resource efficiency through an external judge assessment of deliverable quality and automatically logged operational metrics~\cite{qian2024chatdevcommunicativeagentssoftware, wu2023autogenenablingnextgenllm}.
    \item \textbf{Subjective Team Dynamics:} Evaluating how leadership tone shapes interpersonal agent communication and collaboration climate through sentiment analysis of message logs and post-experiment worker satisfaction surveys~\cite{wageman2005teamdiagnosticsurvey, park2023generativeagentsinteractivesimulacra, li2023camelcommunicativeagentsmind, jiang2024personallminvestigatingabilitylarge, almutairi2025taifaautomatedfeedback}.
\end{itemize}

By comparing these outcomes against a non-leadership Baseline, this research determines whether established human management strategies enhance artificial collaboration or if multi-agent orchestration demands a distinct, machine-centric leadership approach.

\section{Structure}
\label{sec:structure}

The remaining chapters are organized as follows.
Chapter~\ref{chap:background} (Background) establishes the theoretical foundations, synthesizing literature on communicative multi-agent architectures, reviewing the leadership paradigms that lead to Goleman's six-style taxonomy, examining in-context persona enactment in \acp{LLM}, and introducing the measurement instruments.
Chapter~\ref{chap:system-design} (System Design) presents the architecture developed for this study: the four-agent role definitions (Boss, Coder, Writer, and Reviewer), the infrastructure for state management and sandboxed code execution, and the seven-phase collaborative workflow.
Chapter~\ref{chap:experiment-design} (Experiment Design) describes the experimental setup, specifying the 70-run between-subjects matrix ($7 \times 2 \times 5$), the prompt manipulations, the short and long data science tasks, and the embedded data-quality traps.
Chapter~\ref{chap:evaluation} (Evaluation) details the measurement protocols, covering the non-parametric statistical framework, the \ac{LLM}-as-a-judge rubric, operational cost metrics, dual-backend sentiment analysis, the two-stage satisfaction survey, and the validity checks.
Chapter~\ref{chap:results} (Results) reports the findings across all four outcome dimensions.
Chapter~\ref{chap:discussion} (Discussion) interprets them, examines the trade-offs among styles and the sensitivity of the automated evaluators, and states the limitations of the study.
Chapter~\ref{chap:conclusion} (Conclusion) summarizes the contributions, gives architectural recommendations for \ac{MAS} engineering, and outlines directions for future research.

\medskip
\noindent\textbf{Code Availability.}\enspace The complete source code, all seven Boss prompt files, the three worker prompts, the orchestrator, the 70 run artifact folders, and the four analysis notebooks are available in a public repository at \url{https://github.com/helenavioletta/MAS_leadershipstyles}~\cite{mas-leadershipstyles-repo}.